
En
appliquant la formule 2.8, on obtient :
\[
  \frac{\sum_{k\ge0}k\,i_*(\lambda^k(N,y))t^k}
       {1+\sum_{k\ge1} i_*(\lambda^k(N,y))t^k}
  \equiv
  \frac{\sum_{k\ge0}k\lambda^k(i_*(y))t^k}
       {\sum_{k\ge0}\lambda^k(i_*(y))t^k}
       \quad\text{mod.\ torsion}
\]
Si on pose :
\[
  u(t)=1+\sum_{k\ge1}i_*(\lambda^k(N,y))t^k,\qquad
  v(t)=\sum_{k\ge0}\lambda^k(i_*(y))t^k,
\]
cette dernière relation peut encore s'écrire :
\[
  \frac{d}{dt}\!\left(\frac{u(t)}{v(t)}\right)\equiv0
  \quad\text{mod.\ torsion}.
\]
Comme le terme constant de \(u(t)/v(t)\) est \(1\), et qu'on écrit des congruences modulo torsion, on peut intégrer et écrire la congruence \(u(t)\equiv v(t)\). Ceci donne les congruences annoncées.

\textbf{Corollaire 4.4.} \uline{Avec les notations et les hypothèses de} 4.3, \uline{on a pour tout} \(p\ge1\), \uline{et tout} \(y\in K^{\bullet}(Y)\) :
\[
  i_*\bigl(\gamma^p(N,y)\bigr)\equiv \gamma^p(i_*(y))
  \quad\text{mod.\ torsion}.
\]

Le corollaire résulte immédiatement de 4.3, en utilisant l'expression des \(\gamma^p\) comme combinaisons linéaires à coefficients entiers des \(\lambda^p\).

\textbf{Remarque 4.5.} La démonstration même de 4.3 laisse penser que (4.3.1) doit être vrai non seulement modulo torsion, mais aussi en tant qu'égalité. Dans [RRR], l'égalité était montrée en caractéristique \(0\) grâce à une caractérisation géométrique des \(\lambda^p(N,x)\). Une extension de cette méthode en toutes caractéristiques serait sans doute possible si on savait définir les opérations de puissance extérieure sur les complexes parfaits (cf. XIV).

